\chapter*{Preface}
This volume is about the meeting of the integers with probability, and about the work that begins when a theorem is no longer the last word.

The first four parts ask when a deterministic arithmetic object behaves like a random one, and when it does not. Part~I fixes the question: what ``random-like'' should mean, what Cram\'er's model gets right, and how the prime factors of a typical integer become independent. Part~II turns that idea into statistics: random multiplicative functions, primes in short intervals, sieves read as expectations, class groups, and ranks of elliptic curves. Part~III is spectral. The zeros of the zeta function are compared with eigenvalues, the moments of $L$-functions with moments of characteristic polynomials, and $\log|\zeta|$ with a Gaussian. Part~IV is the other direction. Pseudorandomness isolates the structured piece of a function, the Green--Tao theorem puts arbitrarily long progressions inside the primes, and three classical biases show where a random model is subtly wrong.

These thirteen chapters state the questions and the standard theorems that frame them. A statement here is not a proof; where a later chapter uses a result, it says so.

Part~V is a workshop. It follows one problem through its whole life, reads papers with their proofs, tests ideas by computation, compares a heuristic with a theorem, checks a proof, and writes one down. The twelve chapters prove the mathematics they use. The appendices are case studies of conjectures, from twin primes and the Riemann hypothesis to class groups, ranks and the $abc$ conjecture. Appendix~E is read together with the random-matrix arguments of Volume~11, Appendix~B.

The book assumes the series. It names the volume and the chapter where an earlier result is proved, and it does not reprove it. Each workshop chapter opens with the story of its ideas and ends with exercises, a hint, and a solution that explains how to find the argument.
